(1) 半径 \(a\) 的球形各向同性介质（内部介电率 \(\varepsilon_i>0\)、电导率 \(\sigma_i>0\)）埋在无限均匀外介质（\(\varepsilon_o>0,\sigma_o>0\)）中。准静态条件下远场电场在 \(t=0\) 从零阶跃到 \(E_0\hat z\)，两相均无体自由电荷，界面可积累自由面电荷。写出内外偶极形式电势、切向场连续关系以及总法向电流 \((\sigma+\varepsilon\partial_t)E_n\) 的连续方程。
(2) 以 \(K(t)\) 表示外偶极项 \(\phi_o=-E_0r\cos\vartheta+E_0a^3K(t)\cos\vartheta/r^2\)。推导阶跃响应的弛豫时间、初始介电对比极化系数和长时导电对比系数，求诱导偶极矩 \(p(t)\)。
(3) 求界面自由面电荷密度 \(q_s(\vartheta,t)\) 和局部电荷连续性方程；证明初时面电荷为零、长时趋于稳态。对 \(\varepsilon_i=\varepsilon_o,\sigma_i=\sigma_o\) 检验无界面极化极限。
(4) 当 \(K_\varepsilon K_\sigma<0\) 时求唯一偶极反向时间；分析 \(\sigma_i\gg\sigma_o\) 极限，并解释该模型忽略位移电流波传播和界面双电层结构的尺度限制。
